\documentclass[11pt,a4paper]{article}

\usepackage[utf8]{inputenc}
\usepackage[T1]{fontenc}
\usepackage{lmodern}
\usepackage[margin=0.76in]{geometry}
\usepackage[table]{xcolor}
\usepackage{array}
\usepackage{tabularx}
\usepackage{booktabs}
\usepackage{enumitem}
\usepackage{titlesec}
\usepackage{fancyhdr}
\usepackage{lastpage}
\usepackage{hyperref}
\usepackage[normalem]{ulem}

\definecolor{TextInk}{HTML}{233142}
\definecolor{HeadingBlue}{HTML}{1D3557}
\definecolor{AccentTeal}{HTML}{2A7F92}
\definecolor{AccentGold}{HTML}{C58B4E}
\definecolor{RuleGray}{HTML}{D6E0E5}
\definecolor{TableStripe}{HTML}{F3F8FA}
\definecolor{SoftTint}{HTML}{F7F4EE}
\definecolor{SoftBlue}{HTML}{EEF6F8}
\definecolor{SoftGreen}{HTML}{EEF7F0}
\definecolor{SoftRed}{HTML}{FBF0EF}

\hypersetup{
  colorlinks=true,
  linkcolor=HeadingBlue,
  urlcolor=AccentTeal,
  pdftitle={IB Geography SL May 4 Urban And Population Remediation},
  pdfauthor={IB Geography SL Preparation}
}

\color{TextInk}
\setlength{\parskip}{0.45em}
\setlength{\parindent}{0pt}
\setlength{\tabcolsep}{5pt}
\renewcommand{\arraystretch}{1.18}
\setlist[itemize]{leftmargin=1.32em,itemsep=3pt,topsep=4pt}
\setlist[enumerate]{leftmargin=1.45em,itemsep=3pt,topsep=4pt}

\newcolumntype{Y}{>{\raggedright\arraybackslash}X}
\newcolumntype{L}[1]{>{\raggedright\arraybackslash}p{#1}}
\newcolumntype{C}[1]{>{\centering\arraybackslash}p{#1}}

\titleformat{\section}
  {\normalfont\huge\sffamily\bfseries\color{HeadingBlue}}
  {}{0pt}{}
  [\vspace{0.12em}{\color{AccentGold}\titlerule[1pt]}]

\titleformat{\subsection}
  {\normalfont\Large\sffamily\bfseries\color{AccentTeal}}
  {}{0pt}{}

\titleformat{\subsubsection}
  {\normalfont\large\sffamily\bfseries\color{HeadingBlue}}
  {}{0pt}{}

\titlespacing*{\section}{0pt}{1.05em}{0.45em}
\titlespacing*{\subsection}{0pt}{0.78em}{0.22em}
\titlespacing*{\subsubsection}{0pt}{0.55em}{0.16em}

\pagestyle{fancy}
\setlength{\headheight}{14pt}
\fancyhf{}
\fancyhead[L]{\sffamily\color{AccentTeal}Mon 2026-05-04}
\fancyhead[R]{\sffamily\color{HeadingBlue}IB Geography SL}
\fancyfoot[C]{\sffamily\color{HeadingBlue}\thepage\ of \pageref{LastPage}}
\renewcommand{\headrulewidth}{0.4pt}

\newcommand{\term}[1]{\textcolor{HeadingBlue}{\textbf{#1}}}
\newcommand{\exam}[1]{\textcolor{AccentTeal}{\textbf{#1}}}
\newcommand{\todobox}{\fbox{\rule{0pt}{0.8em}\rule{0.8em}{0pt}}}
\newcommand{\boxnote}[3]{%
\noindent\colorbox{#1}{%
  \begin{minipage}{0.965\textwidth}
  \sffamily
  \textbf{\textcolor{HeadingBlue}{#2}} #3
  \end{minipage}
}}

\begin{document}

\begin{center}
  {\sffamily\bfseries\color{HeadingBlue}\LARGE IB Geography SL Study Pack}\\[0.25em]
  {\sffamily\bfseries\color{AccentTeal}\Large Monday 2026-05-04}\\[0.35em]
  {\sffamily\color{TextInk}Urban Environments Remediation + Population Data And Models}
\end{center}

\boxnote{SoftBlue}{How this pack follows the plan.}{The May 4 schedule says:
\term{Audio 1}: weakest \term{Urban environments} content only.
\term{Audio 2}: \term{Population} data, models, and common comparison points.
\term{25 min}: redo the exact question type handled worst on May 2 or May 3.
The May 2 teacher notes identify two Urban priorities: source precision in
short answers and stronger Detroit evidence in the 10-mark answer.}

\section{Today's Target}

\rowcolors{2}{TableStripe}{white}
\begin{tabularx}{\textwidth}{L{0.2\textwidth}Y L{0.25\textwidth}}
\toprule
\textbf{Part} & \textbf{Focus} & \textbf{Output} \\
\midrule
Audio 1 & Fix weakest Urban environments areas: source evidence and
Detroit regeneration / gentrification & Two source-evidence sentences and
three Detroit evidence anchors \\
Audio 2 & Population data, DTM links, population pyramids, migration, and
comparison language & One data-comparison toolkit \\
25 min block & Redo the exact weakest question type from May 2 or May 3 & One
marked rewrite plus an error-log fix \\
\bottomrule
\end{tabularx}
\rowcolors{2}{}{}

\boxnote{SoftRed}{No new-content drift.}{May 4 is a repair day. Do not build a
new set of notes. Use the existing May 1 case-study pack, the May 2 teacher
feedback, and the May 3 error log to make one answer measurably better.}

\section{Audio 1 Notes: Urban Weakness Repair}

\subsection{Weakness 1: Source-Based Precision}

The May 2 teacher feedback awarded the Urban short answers strongly overall,
but Q1 lost a mark because the evidence was too general. For \exam{state} and
\exam{describe} questions, the answer must sound source-based.

\rowcolors{2}{TableStripe}{white}
\begin{tabularx}{\textwidth}{L{0.24\textwidth}Y Y}
\toprule
\textbf{Question type} & \textbf{Risky answer} & \textbf{Better answer} \\
\midrule
\exam{State} evidence & Green space is unequal. & The suburban zone has
\term{34\%} green space, compared with only \term{3\%} in the peripheral
informal settlement. \\
\exam{Describe} contrast & One area has higher rent and another has lower
rent. & The inner-city regeneration zone has a rent index of \term{82},
whereas the peripheral informal settlement has a rent index of \term{24}. \\
\exam{Explain} stress & Rapid growth causes problems. & Rapid in-migration can
raise housing demand faster than formal supply, leading to overcrowding,
informal settlement growth, and service pressure. \\
\bottomrule
\end{tabularx}
\rowcolors{2}{}{}

\subsection{Source Sentence Formula}

Use this formula for short questions:

\begin{center}
{\sffamily\bfseries\color{HeadingBlue}
Place or category + indicator + number or exact source evidence + comparison}
\end{center}

\boxnote{SoftGreen}{Model.}{The peripheral informal settlement has only
\term{3\%} green space and weak service access, while the suburban zone has
\term{34\%} green space. This is direct evidence of uneven environmental
quality.}

\subsection{Weakness 2: Detroit Case-Study Depth}

The May 2 10-mark answer was balanced and relevant, but it needed more precise
case-study evidence. Keep Detroit as the anchor case, but make it more compact
and sharper.

\rowcolors{2}{TableStripe}{white}
\begin{tabularx}{\textwidth}{L{0.18\textwidth}Y}
\toprule
\textbf{Case link} & \textbf{Exam-ready Detroit content} \\
\midrule
Place & Detroit, Michigan, United States \\
Main process & Deindustrialization, population loss, urban decline, selective
regeneration, and gentrification \\
Evidence anchors & Population fell from about \term{1.85 million} in 1950 to
\term{639,111} in the 2020 Census; the May 1 case-study pack also records a
2024 estimate of \term{645,705}, poverty of \term{32.7\%}, and median household
income of about \term{\$39,938} for 2020--2024. \\
Causes & Automobile-industry restructuring, suburbanization, racial and income
segregation, vacancy, fiscal stress, and later selective reinvestment \\
Benefits & Reinvestment in selected districts can improve image, services,
business activity, tax base, reuse of land, and employment opportunities. \\
Limits & Benefits may be spatially uneven; affordability pressure, exclusion,
vacancy, poverty, and weak services can remain outside reinvested districts. \\
\bottomrule
\end{tabularx}
\rowcolors{2}{}{}

\subsection{Detroit Paragraph Upgrade}

\boxnote{SoftTint}{Before.}{Detroit declined and then regeneration helped the
city, but it also displaced people and increased prices.}

\boxnote{SoftGreen}{After.}{Detroit shows that regeneration can partly reduce
urban decline, but it is not a complete social solution. Deindustrialization
and suburbanization helped reduce the city's population from about \term{1.85
million} in 1950 to \term{639,111} in 2020, weakening the tax base and
increasing vacancy. Reinvestment in selected districts can improve image,
services, and economic activity, but the benefits are uneven if lower-income
residents face affordability pressure or if poorer neighborhoods remain
underserved.}

\subsection{Evaluation Sentences To Reuse}

\begin{itemize}
  \item Regeneration is most successful at improving visible economic activity
  in targeted districts, but less successful at solving city-wide inequality.
  \item Gentrification can increase investment and service quality, but it may
  exclude lower-income residents if affordability and participation are weak.
  \item Detroit should be judged as uneven regeneration, not as a simple urban
  success story.
  \item A strategy is more sustainable when physical improvement is matched by
  affordable housing, job access, local participation, and service improvement.
\end{itemize}

\section{Audio 2 Notes: Population Data And Models}

\subsection{Core Terms}

\rowcolors{2}{TableStripe}{white}
\begin{tabularx}{\textwidth}{L{0.24\textwidth}Y}
\toprule
\textbf{Term} & \textbf{Use in answers} \\
\midrule
Population distribution & Where people are located across space. \\
Population density & People per unit area, often people per square kilometre. \\
Fertility rate & Average number of children born per woman. \\
Natural increase & Births minus deaths. \\
Net migration & Immigration minus emigration. \\
Population momentum & Continued growth after fertility falls because a large
young cohort is entering reproductive age. \\
Dependency ratio & Relationship between dependent age groups and the working
age population. \\
\bottomrule
\end{tabularx}
\rowcolors{2}{}{}

\subsection{Data-Table Reading Routine}

Use this four-step routine on Paper 2 population data:

\begin{enumerate}
  \item \term{Extremes}: find highest, lowest, oldest, youngest, fastest
  growth, or fastest decline.
  \item \term{Contrasts}: compare two places using the same indicator.
  \item \term{Model link}: connect to the DTM, population pyramids,
  demographic dividend, ageing, or migration.
  \item \term{Consequence}: explain one likely pressure or opportunity.
\end{enumerate}

\subsection{Demographic Transition Model: Exam Use}

\rowcolors{2}{TableStripe}{white}
\begin{tabularx}{\textwidth}{L{0.18\textwidth}Y Y}
\toprule
\textbf{Stage} & \textbf{Pattern} & \textbf{Likely geography} \\
\midrule
Early transition & High birth rate and falling death rate & Rapid growth,
youthful structure, pressure on services \\
Middle transition & Falling birth rate and lower death rate & Growth slows;
urbanization, education, and contraception may influence fertility \\
Late transition & Low birth and death rates & Stable or slow growth; ageing
begins to matter more \\
Post-transition & Very low fertility and ageing & Possible decline, labour
shortage, pension and healthcare pressure \\
\bottomrule
\end{tabularx}
\rowcolors{2}{}{}

\boxnote{SoftTint}{Model warning.}{The DTM is useful, but it is not automatic.
Migration, policy, conflict, disease, inequality, and culture can alter the
pattern.}

\subsection{Common Comparison Points}

\rowcolors{2}{TableStripe}{white}
\begin{tabularx}{\textwidth}{L{0.26\textwidth}Y Y}
\toprule
\textbf{Comparison} & \textbf{Side A} & \textbf{Side B} \\
\midrule
Age structure & Youthful: high under-15 share and low median age & Ageing:
high 65+ share and high median age \\
Growth driver & Natural increase from high fertility & Net migration or
population momentum despite lower fertility \\
Pressure & Schools, housing, sanitation, maternal health, jobs & Pensions,
healthcare, social care, labour shortage \\
Opportunity & Demographic dividend if jobs and education exist & Skilled
older workforce, service-sector demand, automation incentives \\
Policy response & Education, healthcare, job creation, family planning &
Immigration, later retirement, pension reform, age-friendly planning \\
\bottomrule
\end{tabularx}
\rowcolors{2}{}{}

\subsection{Model Population Table}

The data below are illustrative and designed for practice.

\rowcolors{2}{TableStripe}{white}
\begin{tabularx}{\textwidth}{L{0.24\textwidth}C{0.17\textwidth}C{0.17\textwidth}C{0.17\textwidth}}
\toprule
\textbf{Indicator} & \textbf{Country Alpha} & \textbf{Country Beta} & \textbf{Country Gamma} \\
\midrule
Total fertility rate & 4.6 & 1.3 & 2.1 \\
Median age & 18 & 48 & 32 \\
Annual population growth & 2.7\% & -0.3\% & 0.9\% \\
Net migration rate & -1.2\% & 0.4\% & 0.8\% \\
Population aged under 15 & 43\% & 12\% & 24\% \\
Population aged 65 and over & 3\% & 29\% & 12\% \\
\bottomrule
\end{tabularx}
\rowcolors{2}{}{}

\boxnote{SoftGreen}{Model comparison.}{Country Alpha is much more youthful
than Country Beta: Alpha has a median age of \term{18} and \term{43\%} under
15, whereas Beta has a median age of \term{48} and \term{29\%} aged 65 and
over. Alpha may face pressure on schools and future employment, while Beta may
face pension, healthcare, and labour-supply pressure.}

\section{25-Minute Remediation Block}

\rowcolors{2}{TableStripe}{white}
\begin{tabularx}{\textwidth}{L{0.16\textwidth}Y L{0.22\textwidth}}
\toprule
\textbf{Time} & \textbf{Task} & \textbf{Output} \\
\midrule
0--3 min & Re-read the May 2 teacher feedback and choose the exact weakness:
source precision, Detroit evidence, or May 3 Paper 2 issue. & One chosen
target \\
3--19 min & Redo \term{one} selected workbook task. For Task B, this includes
the \term{4 minute} plan and \term{12 minute} write. & One timed rewrite \\
19--23 min & Mark it using the May 4 marking guide. & Score and reason \\
23--25 min & Write a one-line fix for the final error log. & One sentence to
reuse on exam day \\
\bottomrule
\end{tabularx}
\rowcolors{2}{}{}

\boxnote{SoftTint}{Strict timing rule.}{Inside the \term{25 minute} block,
complete \term{one} primary redo: Task A, Task B, or Task C. If Task A is
chosen and finished early, use the spare minutes to write two alternate
source-evidence sentences from the same table; do not start a second full task
inside the timed block.}

\section{Sources And Carry-Forward}

\begin{itemize}
  \item Local plan: the April 23 to May 11 preparation plan in the repo root.
  \item Local feedback: the May 2 Paper 1 simulation teacher notes in
  \texttt{plans}.
  \item Local Detroit source list and case-study anchors: the May 1 case-study
  bank study pack in \texttt{plans}.
  \item After this session, carry only the corrected answer pattern forward:
  source evidence for short questions, named Detroit evidence for 10-mark
  answers, and one Population data-comparison routine.
\end{itemize}

\end{document}
